\documentclass[a4paper,12pt,titlepage,finall]{article}

\usepackage[T1,T2A]{fontenc}     % форматы шрифтов
\usepackage[utf8x]{inputenc}     % кодировка символов, используемая в данном файле
\usepackage[russian]{babel}      % пакет русификации
\usepackage{tikz}                % для создания иллюстраций
\usepackage{pgfplots}            % для вывода графиков функций
\usepackage{geometry}		 % для настройки размера полей
\usepackage{indentfirst}         % для отступа в первом абзаце секции
\usepackage{minted}

% выбираем размер листа А4, все поля ставим по 3см
\geometry{a4paper,left=30mm,top=30mm,bottom=30mm,right=30mm}


\usepgfplotslibrary{fillbetween} % для изображения областей на графиках

\begin{document}
% Титульный лист
\begin{titlepage}
    \begin{center}
	{\small \sc Московский государственный университет \\имени М.~В.~Ломоносова\\
	Факультет вычислительной математики и кибернетики\\}
	\vfill 
	{\Large \sc Отчет по заданию №6}\\
	~\\
	{\large \bf <<Сборка многомодульных программ. \\
	Вычисление корней уравнений и определенных интегралов.>>}\\ 
	~\\
	{\large \bf Вариант 6 / 4 / 3}
    \end{center}
    \begin{flushright}
	\vfill {Выполнил:\\
	студент 103 группы\\
	Огай~В.~Л.\\
	~\\
	Преподаватель:\\
	Кузьменкова~Е.~А.\\}
    \end{flushright}
    \begin{center}
	\vfill
	{\small Москва\\2024}
    \end{center}
\end{titlepage}

% Автоматически генерируем оглавление на отдельной странице
\tableofcontents
\newpage

\section{Постановка задачи}

С заданной точностью {$\varepsilon$} вычислить площадь плоской фигуры, ограниченной тремя кривыми: 
\begin{itemize}
\item[-]{$f_1(x) = 0.6{x} + 3$}
\item[-]{$f_2(x) = ({x} - 2)^3 - 1$}
\item[-]{$f_3(x) = \frac{3}{x}$}
\end{itemize}
Для вычисления площади плоской фигуры необходимо реализовать численные методы, предусмотренные вариантом задания. В данном варианте $-$ комбинированный метод (хорд и касательных) для определения корней уравнений и метод парабол (метод Симпсона) для вычисления определенного интеграла от функции на отрезке.
Отрезок  {$[a,b]$}, содержащий корень, должен быть найден вручную с необходимым математическим обоснованием. Площадь фигуры необходимо вычислить с точностью {$\varepsilon = 0.001$}. {$\varepsilon_1$} $-$ точность, с которой считаются корни, и {$\varepsilon_2$} $-$ погрешность вычисления определенного интеграла, также должны быть выбраны с соответствующим математическим обоснованием.



\newpage

\section{Математическое обоснование}
\subsection{Отрезки для корней и описание методов}
В варианте 6 представленного задания задана область, изображенная на (Рис.1), площадь которой необходимо найти. Для этого требуется вычислить абсциссы точек пересечения кривых, которые являются решением уравнения {$F(x) = f(x) - g(x) = 0$}. Для поиска корней данного уравнения используется комбинированный метод (хорд и касательных). \\\par

\begin{figure}[h]
\centering
\begin{tikzpicture}
\begin{axis}[% grid=both,                % рисуем координатную сетку (если нужно)
              axis lines=middle,          % рисуем оси координат в привычном для математики месте
             restrict x to domain=-10:10,  % задаем диапазон значений переменной x
             restrict y to domain=0:10,  % задаем диапазон значений функции y(x)
             axis equal,
             enlargelimits,              % разрешаем при необходимости увеличивать диапазоны переменных
             legend cell align=left,     % задаем выравнивание в рамке обозначени
             xmin = -5,
             xmax = 5,
             ymin = 0,
             scale=2,                    % задаем масштаб 2:1
             ]                   

% первая функция
% параметр samples отвечает за качество прорисовки
\addplot[green,domain=-3:9,samples=256,thick] {0.6*x + 3};
% описание первой функции
\addlegendentry{$f_1=0.6{x} + 3$}

\addlegendimage{empty legend}\addlegendentry{}

\addplot[blue,domain=-3:9,samples=256,thick] {(x-2)^3 - 1};
\addlegendentry{$f_2=(x-2)^3 - 1$}

\addplot[red,domain=-3:9,samples=256,thick] {3/x};
\addlegendentry{$f_3=\frac{3}{x}$}

\addlegendimage{empty legend}\addlegendentry{}

\end{axis}
\end{tikzpicture}
\caption{Плоская фигура, ограниченная графиками заданных уравнений}
\label{plot1}
\end{figure}

Стоит различать два случая:\par
1. Первая и вторая производные функции {$F(x)$} имеют одинаковый знак.\par
2. Первая и вторая производные функции {$F(x)$} имеют разные знаки.\\\par

Комбинированный метод (хорд и касательных) $-$ итеративный метод приближения корня уравнения {$F(x) = 0$} на отрезке {$[a, b]$}. Требуется, чтобы функция была непрерывна на {$[a, b]$}, а вторая производная функции была монотонна на {$[a, b]$}. В случае 1 полагается, что 
\begin{center}
    {$a_{n+1} = a_n - \frac{f(a_n)(b_n-a_n)}{f(b_n)-f(a_n)} $,
    {$b_{n+1} = b_n - \frac{f(b_n)}{f^`(b_n)}$}.
\end{center}
В случае 2: 
\begin{center}
    {$a_{n+1} = a_n - \frac{f(a_n)}{f^`(a_n)}$},
    {$b_{n+1} = b_n - \frac{f(b_n)(b_n-a_n)}{f(b_n)-f(a_n)}$}.
\end{center}\par

Критерием остановки является условие {$|b_{n+1} - a_{n+1}| < \varepsilon_1$}. \cite{prog}\\\par

Здесь и далее \begin{itemize}
    \item[-] {$F_1(x) = f_1(x) - f_2(x)$;}
    \item[-] {$F_2(x) = f_1(x) - f_3(x)$;}
    \item[-] {$F_3(x) = f_2(x) - f_3(x)$.}
\end{itemize}
Для поиска корней были выбраны отрезки: для {$F_1(x)$} $-$ {$[2.5, 4]$}, для {$F_2(x)$} $-$ {$[0.5, 2.5]$}, для {$F_3(x)$} $-$ {$[2.5, 4]$}. Проверим, что все функции удовлетворяют на нем необходимым условиям, a именно: на концах отрезка функция \textit{$F(x) = f(x) - g(x) $} принимает значения разных знаков, ее первые и вторые производные сохраняют знак.\\\par

\begin{enumerate}
\item {$F_1(x) = 0.6x + 3 - ((x-2)^3 - 1)$}.
\begin{enumerate}
\item {$F_1(2.5) = 5.375 > 0, F_1(4) = -1.6 < 0$}
\item {$F_1'(x) = 0.6 - 3(x-2)^2 < 0$} {$\forall{x} \in [2.5, 4]$}
\item {$F_1''(x) = 6(x - 2) > 0 $} {$\forall{x} \in [2.5, 4] $}
\end{enumerate}\par

\item {$F_2(x) = 0.6x + 3 - \frac{3}{x}$}.

\begin{enumerate}
\item {$F_2(0.5) = -2.7 < 0, F_2(2.5) = 3.3 > 0$}
\item {$F_2'(x) = 0.6 + \frac{3}{x^2} > 0$} {$\forall{x} \in [0.5, 2.5]$}
\item {$F_2''(x) = -\frac{6}{x^3} < 0 $} {$\forall{x} \in [0.5, 2.5] $}
\end{enumerate}\par

\item {$F_3(x) = (x-2)^3 - 1 - \frac{3}{x}$}.
\begin{enumerate}
\item {$F_3(2.5) = -2.075 < 0, F_1(4) = 6.25 > 0$}
\item {$F_3'(x) = 3(x-2)^2 + \frac{3}{x^2} > 0$} {$\forall{x} \in [2.5, 4]$}
\item {$F_3''(x) = 6(x-2) - \frac{6}{x^3} < 0 $} {$\forall{x} \in [2.5, 4]$}
\end{enumerate}\par
\end{enumerate}\par
Далее, после того как были найдены абсциссы точек пересечения кривых, необходимо вычислить определенные интегралы, используя метод парабол (метод Симпсона).\par

Метод парабол (метод Симпсона) позволяет приближённо вычислить определённый интеграл, прибегая к аппроксимации частей кривой многочленами второй степени.\cite{math} На отрезке {$[a, b]$} формула для определенного интеграла {$I$} следующая: \\\par
\begin{center}
    {$I_{2n} = \frac{h}{3}(f(a) + f(b)) + 4(f(x_1) + f(x_3) +... f(x_{n - 1})) + 2(f(x_2) + f(x_4) ... + f(x_{n-2}))$}, где {$x_i = a + ih$}, {$h = (b - a) / n$}, $n$ - чётно.
\end{center}
Критерий остановки в данном случае - правило Рунге:
\begin{center}
    $|I - I_n| \approx \theta|I_{2n} - I_n|$,
\end{center}
где $\theta$ равно $\frac{1}{15}$ для данного метода. Критерием остановки является условие $\theta|I_n - I_{2n}|$ < $\varepsilon_2$. В таком случае метод возвращает $I_{2n}$ в качестве ответа. \cite{prog}
\subsection{Погрешности}
Аддитивность абсолютных погрешностей при суммировании измеряемых результатов известна. Следовательно, ошибка при вычислении интегралов будет равна $3\varepsilon_2$. Выберем $\varepsilon_1$. Интегрируя функцию $F$ в пределах $x_1, x_2$, которые были получены функцией \textit{root}, получаем, что ошибка при вычислении $\leq$ суммы $\varepsilon_1 * max(F(x_1 \pm \varepsilon_1))$ и $\varepsilon_1 * max(F(x_2 \pm \varepsilon_1))$ $-$ площадей прямоугольников (это значит, что избыточные или недостающие криволинейные трапеции будут лежать внутри этих прямоугольников). Нетрудно заметить, что на выбранном участке график $f_1$ лежит не ниже остальных графиков, следовательно погрешностями при интегрировании $f_1$ можно оценить сверху погрешности при интегрировании остальных функций. В силу достаточной малости ошибочного интервала и непрерывности $f_1$, любое используемое значение функции $f_1$ можно оценить сверху константой $a$ такой, что прямая $y = a$ будет лежать сильно выше рассматриваемой фигуры. В данном случае, как видно из графика и приблизительных значений ординат точек пересечения графиков (п. 3), достаточно взять $a = 10$. Таким образом, мы можем рассмотреть сумму площадей 6 прямоугольников со сторонами $\varepsilon_1$ и $a$ (их шесть, поскольку учитываются 3 интеграла). Тогда общая ошибка не будет превосходить
\begin{center}
    $60\varepsilon_1 + 3\varepsilon_2$.
\end{center}
Возьмем $\varepsilon_1 = \varepsilon_2 = 10^{-5}$.
Подставив эти значения в полученную выше оценку, мы получаем $63 \cdot 10^{-5} < 10^{-3}$, следовательно подобранные значения обеспечат требуемую точность.
\newpage

\section{Результаты экспериментов}

Результаты эксперимента представлены в таблице. 

\begin{table}[h]
\centering
\begin{tabular}{|c|c|c|}
\hline
Кривые & $x$ & $y$ \\
\hline
1 и 2 &  3.847760 & 5.308656 \\
1 и 3 &  0.854103 & 3.512462 \\
2 и 3 &  3.243928 & 0.924805 \\
\hline
\end{tabular}
\caption{Координаты точек пересечения}
\label{table1}
\end{table}

\begin{figure}[h]
\centering
\begin{tikzpicture}
\begin{axis}[% grid=both,                % рисуем координатную сетку (если нужно)
              axis lines=middle,          % рисуем оси координат в привычном для математики месте
             restrict x to domain=-10:10,  % задаем диапазон значений переменной x
             restrict y to domain=0:10,  % задаем диапазон значений функции y(x)
             axis equal,
             enlargelimits,              % разрешаем при необходимости увеличивать диапазоны переменных
             legend cell align=left,     % задаем выравнивание в рамке обозначени
             xmin = -5,
             xmax = 5,
             ymin = 0,
             scale=2,                    % задаем масштаб 2:1
             xticklabels={,,},           % убираем нумерацию с оси x
             yticklabels={,,}]           % убираем нумерацию с оси y
% первая функция
% параметр samples отвечает за качество прорисовки
\addplot[green,domain=-5:5,samples=256,thick,name path=A] {0.6*x + 3};
% описание первой функции
\addlegendentry{$f_1=0.6{x} + 3$}

% добавим немного пустого места между описанием первой и второй функций


% вторая функция
% здесь необходимо дополнительно ограничить диапазон значений переменной x
\addplot[blue,domain=-5:5,samples=256,thick,name path=B] {(x-2)^3- 1};
\addlegendentry{$f_2=(x-2)^3 - 1$}

% дополнительное пустое место не требуется, так как формулы имеют небольшой размер по высоте

% третья функция
\addplot[red,domain=-5:5,samples=256,thick,name path=C] {3/x};
\addlegendentry{$f_3=\frac{3}{x}$}
\addlegendimage{empty legend}\addlegendentry{}
% закрашиваем фигуру
\addplot[blue!20,samples=256] fill between[of=A and C,soft clip={domain=0.854103:3.243928}];
\addplot[blue!20,samples=256] fill between[of=A and B,soft clip={domain=3.243928:3.847760}];
\addlegendentry{$S=7.488420$}

% Поскольку автоматическое вычисление точек пересечения кривых в TiKZ реализовать сложно,
% будем явно задавать координаты.
\addplot[dashed] coordinates { (0.854103, 3.512462) (0.854103, 0) };
\addplot[color=black] coordinates {(0.854103, 0)} node [label={-90:{\small 0.854}}]{};

\addplot[dashed] coordinates { (3.243928, 0.924805) (3.243928, 0) };
\addplot[color=black] coordinates {(3.243928, 0)} node [label={-90:{\small 3.24}}]{};

\addplot[dashed] coordinates { (3.847760, 5.308656) (3.847760, 0) };
\addplot[color=black] coordinates {(3.847760, 0)} node [label={-90:{\small 3.848}}]{};

\end{axis}
\end{tikzpicture}
\caption{Плоская фигура, ограниченная графиками заданных уравнений}
\label{plot2}
\end{figure}

\newpage

\section{Структура программы и спецификация функций}

Разбиение программы на модули и функции, и связи между ними приведены на рис.~\ref{plot3}.
\begin{figure}[h]
\centering
\begin{tikzpicture}[node distance={30mm}, main/.style = {draw, rectangle}] 
\node[main] (1) {integral.c}; 
\node[main] (2) [right of=1] {functions.asm};
\node[main] (3) [above right of=1] {Makefile};
\node[main] (4) [above left of=1] {integral, root}; 
\node[main] (5) [below of=2] {$f_{1,2,3}$; $f'_{1, 2, 3}$};

\draw[->] (1) -- (3);
\draw[<-] (1) -- (4);
\draw[->] (2) -- (3);
\draw[<-] (2) -- (5);
\end{tikzpicture}

\caption{Разбиение программы на компоненты}
\label{plot3}
\end{figure}

Файлы и используемые в них функции:

\begin{minted}[frame=lines]{python}
integral.c:

g2(double x), g3(double x)  - возвращают значения g2 и g3
dg1(double x), dg2(double x), dg3(double x)
                            - возвращают значения производных от g1, g2 и g3
root(double f(double), double g(double), double df(double), 
double dg(double), double a, double b, double eps1)        
                            - возвращает корень f(x)-g(x)=0
integral(double f(double), double a,double b, double eps2)       
                            - возвращает интеграл, с заданной точностью
help(void)                  - выводит информацию о программе
main(int argc, char **argv) - считает необходимые корни и интегралы, 
                              выводит ответ к задаче, 
                              разбирает опции командной строки
\end{minted}

\begin{minted}[frame=lines]{python}
functions.asm:

f1, f2, f3, g1      - возвращают f1(x), f2(x), f3(x), g1(x)
df1, df2, df3       - возвращают производные от f1(x), f2(x), f3(x)
\end{minted}
Программой поддерживаются следующие опции командной строки (и их короткие вариации):
\begin{itemize}
    \item \verb|--root, -r| - вывод абсцисс точек пересечения данных функций;
    \item \verb|--iterations, -i| - вывод количества итераций, потребовавшегося для вычисления корней уравнения;
    \item \verb|--help, -h| - вывод информации о программе;
    \item \verb|--test-root, -R| - опция, принимающая аргумент вида \verb|F1:F2:A:B:E:R| и тестирующая функцию \textit{root} с помощью данных из этого аргумента;
    \item \verb|--test-integral, -I| - опция, принимающая аргумент вида \verb|F:A:B:E:R| и тестирующая функцию \textit{integral} с помощью даннымх из этого аргумента.
\end{itemize}
Последние две опции выводят результат работы тестируемого метода, абсолютную и относительную погрешности. Программа, запущенная без опций, выводит ответ на задачу (площадь плоской фигуры).
\newpage
\section{Сборка программы (Make-файл)}
Код Makefile (за исключением предопределённых в задании флагов):

\begin{minted}[frame=lines]{make}
all: integral

integral: integral.o functions.o
	gcc -m32 -lm integral.o functions.o -o integral
	
integral.o: 
	gcc -m32 -c integral.c -o integral.o
	
functions.o: 
	nasm -f elf32 functions.asm -o functions.o

test: integral
	./$< -R 1:2:0.2984:3.9999:0.00001:0.697082
	./$< -R 1:3:0.01:0.3308:0.00001:0.330055
	./$< -R 2:3:0.666:7.777:0.00001:0.922758
	./$< -I 1:4.5:5.5:0.0001:-370.752
	./$< -I 2:2.37:7.943:0.0001:11143.2
	./$< -I 3:5.6728:1.5684:0.0001:-0.0576641

clean:
	rm *.o integral
\end{minted}
	
Цель \textit{all} зависит от исполняемого файла \textit{integral}, который, в свою очередь, зависит от объектных файлов \textit{integral.o} и \textit{functions.o}. Эти объектные файлы зависят от \textit{integral.c} и \textit{functions.asm} соответственно. Данная цель собирает исполняемый файл\par
Цель \textit{test} зависит от этого же исполняемого файла и запускает его с ключами test-root и test-integral (короткие вариации). В случае отсутствия исполняемого файла цель \textit{test} пройдёт по реквизитам \textit{integral} и соберёт исполняемый файл, после чего запустит этот файл с указанными выше ключами\par
Цель \textit{clean} очищает все промежуточные файлы.\par
При сборке исполняемого файла с помощью \textit{Makefile}, сначала проводится поиск файлов с полным совпадением имён, указанных в целях, реквизитах и командах. При их наличии новый файл не будет создан, следовательно при изменении исходных файлов (на Си или Асме) стоит поочерёдно воспользоваться целями \textit{clean} и \textit{all} перед работой с исполняемым файлом.\par
Также стоит учесть необходимость использования зависимостей (реквизитов) в \textit{integral}, т.к. каждая цель, использующая \textit{integral}, для корректной работы должна собирать исполняемый файл \textit{integral} в случае его отсутствия.
\newpage

\section{Отладка программы, тестирование функций}

Функция \textit{root} тестировалась на отдельных функциях, на произвольных отрезках, удовлетворяющих условиям применения комбинированного метода (хорд и касательных). Обоснование корректности выбранных отрезков было приведено в п. 2.1. Новые тест-функции и их производные, а также результаты работы программы и результаты, полученные с помощью сервиса WolframAlpha приведены далее.

Для поиска корней были выбраны отрезки: для {$g_1(x)$} $-$ {$[0.2984, 3.9999]$}, для {$g_2(x)$} $-$ {$[0.01, 0.3308]$}, для {$g_3(x)$} $-$ {$[0.666, 7.777]$}. \\\par

\begin{enumerate}
\item {$g_1(x) = ln(x) - 15x^2 + 3.89$}.
\item {$g_2(x) = 13x^3 - 12x^2 + 11x - 10$}.
\item {$g_3(x) = \frac{1}{8x^2}$}.
\end{enumerate}\par

\begin{enumerate}
\item {$g'_1(x) = \frac{1}{x} - 30x$}.
\item {$g'_2(x) = 39x^2 - 24x + 11$}.
\item {$g'_3(x) = -\frac{1}{4x^3}$}.
\end{enumerate}\par

\begin{table}[h]
\centering
\begin{tabular}{|c|c|c|}
\hline
Номера функций & \textit{root} & WolframAlpha \\
\hline
    1-2 & 0.697082 & 0.697082 \\
    1-3 & 0.330056 & 0.330056 \\
    2-3 & 0.922758 & 0.922758 \\
\hline
\end{tabular}
\caption{результаты тестирования функции \textit{root}}
\label{table2}
\end{table}

Функция \textit{integral} также тестировалась на данных функциях, также с произвольными пределами интегрирования: \par
\begin{table}[h]
\centering
\begin{tabular}{|c|c|c|c|}
\hline
Функция & Отрезок интегрирования & integral & Wolfram Alpha \\
\hline
{$f_1(x)$} & {$[4.5, 5.5]$} & -370.752234 & -370.752 \\
{$f_2(x)$} & {$[2.37, 7.943]$} & 11143.199406 & 11143.2 \\
{$f_3(x)$} & {$[5.6728, 1.5684]$} & -0.057755 & -0.0576641 \\

\hline
\end{tabular}
\caption{результаты тестирования функции \textit{integral}}
\label{table1}
\end{table}

Численные методы и функции были отлажены утилитой gdb, встроенной в SASM и CLion.

\newpage

\section{Анализ допущенных ошибок}
В процессе выполнения работы были допущены следующие ошибки:
\begin{enumerate}
\item {В \textit{Makefile integral} не имел реквизитов, из-за чего при запуске цели \textit{test} без предварительного запуска цели \textit{all} программа не запускалась}.

\end{enumerate}

\newpage
\begin{raggedright}
\addcontentsline{toc}{section}{Список цитируемой литературы}
\begin{thebibliography}{99}
\bibitem{prog} Трифонов~Н.\,А., Пильщиков~В.\,Н. Задания практикума на ЭВМ (1 курс). Учебное пособие, 2-е исправленное издание. —
М.: МГУ, 2001. — 32 с. 
\bibitem{math} Ильин~В.\,А., Садовничий~В.\,А., Сендов~Бл.\,Х. Математический анализ. Т.\,1~---
    Москва: Наука, 1985.
\end{thebibliography}
\end{raggedright}

\end{document}